% Figure: coherent type-1 feed-forward loop. X activates Y; X and Y
% both activate Z; the loop filters brief excursions of X.
\begin{figure}[htbp]
\centering
\begin{tikzpicture}[
  node distance=2.6cm,
  reg/.style={circle, draw=engblack, very thick, fill=engblue!15,
    minimum size=1.0cm, font=\bfseries},
  act/.style={-{Latex[length=2.2mm]}, thick, draw=engblack},
  lbl/.style={font=\scriptsize, align=left},
]

% Topology (left)
\begin{scope}[shift={(0,0)}]
\node[reg] (x) at (0,2.6) {X};
\node[reg] (y) at (-1.4, 0.6) {Y};
\node[reg] (z) at (1.4, 0.6) {Z};
\draw[act] (x) -- (y);
\draw[act] (x) -- (z);
\draw[act] (y) -- (z);
\node[font=\small\bfseries] at (0, -0.6) {Topology};
\end{scope}

% Step-up response (middle)
\begin{scope}[shift={(5.5,0)}]
\draw[->, thick] (-0.2, 0) -- (3.6, 0);
\draw[->, thick] (-0.2, 0) -- (-0.2, 3.0);
% X step-up
\draw[engblue, very thick] (-0.2, 0.4) -- (0.6, 0.4)
  -- (0.6, 2.5) -- (3.4, 2.5);
\node[font=\scriptsize, engblue, anchor=west] at (3.5, 2.5) {X};
% Z delayed
\draw[engred, very thick]
  (-0.2, 0.4) -- (0.6, 0.4)
  to[out=15, in=180] (1.5, 0.5)
  to[out=0, in=200] (2.5, 2.0)
  to[out=20, in=180] (3.4, 2.3);
\node[font=\scriptsize, engred, anchor=west] at (3.5, 2.3) {Z};
\node[font=\scriptsize, anchor=south] at (1.5, 2.7) {delay};
\node[font=\small\bfseries] at (1.5, -0.6) {Step-up};
\node[font=\scriptsize, anchor=west] at (-0.2, 3.2) {amplitude};
\node[font=\scriptsize, anchor=west] at (3.6, -0.05) {time};
\end{scope}

% Step-down response (right)
\begin{scope}[shift={(10.5,0)}]
\draw[->, thick] (-0.2, 0) -- (3.6, 0);
\draw[->, thick] (-0.2, 0) -- (-0.2, 3.0);
% X step-down
\draw[engblue, very thick] (-0.2, 2.5) -- (1.0, 2.5)
  -- (1.0, 0.4) -- (3.4, 0.4);
\node[font=\scriptsize, engblue, anchor=west] at (3.5, 0.4) {X};
% Z drops promptly
\draw[engred, very thick]
  (-0.2, 2.3) -- (1.0, 2.3)
  to[out=-30, in=180] (1.4, 0.5) -- (3.4, 0.5);
\node[font=\scriptsize, engred, anchor=west] at (3.5, 0.5) {Z};
\node[font=\scriptsize, anchor=south] at (1.4, 1.5) {prompt};
\node[font=\small\bfseries] at (1.5, -0.6) {Step-down};
\node[font=\scriptsize, anchor=west] at (3.6, -0.05) {time};
\end{scope}

\end{tikzpicture}
\caption[Coherent type-1 feed-forward loop]{The coherent type-1
feed-forward loop is the most common motif in the \emph{E.\ coli}
transcription network \cite{paper:v6c01-shoval-alon-2010}. The
regulator $X$ activates both $Y$ and the output $Z$, and $Y$ in turn
activates $Z$ (left). The loop responds asymmetrically: a step-up of
$X$ produces a \emph{delayed} activation of $Z$ (middle), because
$Z$ requires both $X$ and $Y$ and $Y$ takes time to accumulate; a
step-down of $X$ produces a \emph{prompt} deactivation of $Z$ (right).
The asymmetry implements a sign-sensitive delay that filters brief
excursions of $X$ without filtering brief drops.}
\label{fig:vol06:ch01:feedforward-loop}
\end{figure}
